\documentclass[12pt]{article}

\usepackage[a4paper, margin=2.5cm]{geometry}
\usepackage{amsmath}
\usepackage{amssymb}
\usepackage{amsthm}
\usepackage[colorlinks=true, linkcolor=blue, citecolor=blue, urlcolor=blue]{hyperref}
\usepackage{booktabs}
\usepackage{enumitem}
\usepackage{natbib}

\newtheorem{theorem}{Theorem}
\newtheorem{proposition}[theorem]{Proposition}
\newtheorem{corollary}[theorem]{Corollary}
\newtheorem{definition}[theorem]{Definition}
\newtheorem{remark}[theorem]{Remark}
\newtheorem{hypothesis}[theorem]{Hypothesis}

\newcommand{\Fcal}{\mathcal{F}}
\newcommand{\Ncl}{\mathcal{N}_{\mathrm{cl}}}
\newcommand{\VW}{V_{\mathrm{W}}}
\newcommand{\HW}{H_{\mathrm{W}}}
\newcommand{\Sadm}{\mathcal{S}_{\mathrm{adm}}}
\newcommand{\Rst}{R_{\mathrm{st}}}
\newcommand{\dUrel}{\delta U_{\mathrm{rel}}}
\newcommand{\Mcl}{\mathcal{M}_{\mathrm{cl}}}
\newcommand{\Ocal}{\mathcal{O}}
\newcommand{\twoI}{2I}
\newcommand{\sigmaG}{\sigma_{\mathrm{Gal}}}

\title{From Closure Dynamics to Quantum Structure:\\[6pt]
\large A Structural Synthesis of the Dissipative-to-Schr\"odinger Programme}
\author{%
  Lee Smart\\[4pt]
  \textit{Institute of Vibrational Field Dynamics}\\[2pt]
  \texttt{contact@vibrationalfielddynamics.org}\\[2pt]
  \texttt{@vfd\_org}%
}
\date{April 2026}

\begin{document}

\maketitle

\noindent\textbf{Status:} Preprint (not peer-reviewed)

\begin{abstract}
\textbf{Recovery position.} The Schr\"odinger equation is the correct description of non-relativistic quantum dynamics within its domain of validity. This paper does not replace it, modify it within that domain, or propose an alternative. The question addressed is structural: \emph{why does the Schr\"odinger equation take the form it does, and from what deeper object is it inherited?} Within the VFD cascade $E_8 \to H_4 \to 40 \to D_4 \to 16 \to 8 \to 0$ on the 600-cell polytope (Paper~XXXVI Foundations F1--F8), two conditional routes yield Schr\"odinger-type evolution under different hypotheses.

\textbf{Route A (stochastic, Papers~XVI--XX).} Starting from a dissipative gradient-flow Langevin system, phase complexification, a Nelson-type forward/backward pairing, and an equilibrium-centred normal form, one obtains an exact nonlinear wave equation whose equilibrium tangent reproduces $i\,\partial_t\Psi = \HW\Psi$ with the Witten Hamiltonian $\HW = -(\sigma^2/2)\Delta + \VW$.

\textbf{Route B (cascade-structural, conditional).} If one adopts as a modelling postulate that time evolution on the H${}_4$ rung is generated by the H${}_4$ graph Laplacian $\Delta_{H_4}$ on $\mathcal{H} = L^2(\text{600-cell vertices})$, then $\Delta_{H_4}$ is self-adjoint (trivially, by real-symmetry), and Stone's theorem on one-parameter unitary groups yields a Schr\"odinger-type evolution $i\hbar\,\partial_t\psi = \hat H\psi$ with $\hat H = \hbar\,\Delta_{H_4}$. Paper~XXXVI Foundation~F5 ($\sigmaG$-intertwining of the scalar closure functional, where $\sigmaG$ is the Galois/lattice automorphism; see the Notation note below) supplies a discrete symmetry of the closure functional under Paper~XXXVI's hypotheses on the coefficient structure (rational coefficients, $\mathbb{Q}(\varphi)$-defined ambient data); any induced symmetry of the chosen generator is conjectural here, not proved, and is not used to produce the unitary group, which rests on self-adjointness of the chosen generator.

The two routes are \emph{parallel structural paths}, not competing derivations of the same equation. Route~A identifies $\hat H$ with the Witten Hamiltonian $\HW$ from closure dynamics; Route~B posits $\hat H = \hbar\,\Delta_{H_4}$ from the H${}_4$ graph geometry. Reconciliation of the two Hamiltonians at the level of effective parameters is an open question (Section~\ref{sec:open}). What the routes share, in their stated linear regimes, is the generic form of the equation: Route~B has a self-adjoint generator outright, while Route~A does so only under the equilibrium-tangent and finite-dimensional invariant-subspace hypotheses of Section~\ref{sec:geometric}.

A geometric-quantum section (Section~\ref{sec:geometric}) makes explicit the K\"ahler structure at the equilibrium tangent \emph{under a sector-subspace invariance hypothesis}: the projective state manifold is $\mathbb{CP}^{N-1}$, the linearised Schr\"odinger flow is a Hamiltonian flow on this manifold with respect to the imaginary-part-of-inner-product symplectic form, and the induced projective action of $U(N)$ factors through $\mathrm{PU}(N)$ and acts by K\"ahler isometries. This is standard geometric quantum mechanics \citep{AshtekarSchilling1999,KobayashiNomizu1969} applied to the invariant finite-dimensional subspace; the closure-derived content is the identification of $\HW$ as the effective Hamiltonian on that subspace. Paper~XXX uses this equilibrium-tangent structure as an input to its K\"ahler-uniqueness argument for the Born rule via Petz's theorem. The full nonlinear dynamics of Route~A remains non-unitary, as previously established; Route~B's unitarity is the cascade-projective statement valid at the H${}_4$ rung.
\end{abstract}

%==================================================================
\section{Introduction}\label{sec:intro}
%==================================================================

The Schr\"odinger equation occupies a central role in modern physics. It is typically introduced as a fundamental postulate governing wavefunction evolution. A long-standing question in mathematical physics is whether such dynamics can arise from deeper structural principles --- either as an equilibrium-tangent recovery from a broader dynamical system (Route~A, pursued in Papers~XII--XX) or as an operator-theoretic consequence of a separately posited self-adjoint generator on a chosen Hilbert space (Route~B) --- rather than being taken as the starting dynamical law.

The closure programme (Papers~XII--XX) approaches this question from a specific direction: starting from dissipative gradient-flow systems and stochastic dynamics, and asking what happens when the resulting stochastic structure is taken seriously as a paired forward/backward system.

The purpose of this paper is to synthesise the technical results of Papers~XVI--XX~\citep{SmartXVI, SmartXVII, SmartXVIII, SmartXIX, SmartXX} into a coherent structural statement, identify the precise relationship to known quantum formulations, and clearly delineate what has and has not been achieved.

\subsection*{Core question}

\begin{quote}
Can Schr\"odinger-type dynamics be recovered from deeper structure, either as an equilibrium-tangent limit of a broader (non-unitary) dynamical system, or as an operator-theoretic consequence of a separately posited self-adjoint generator on a chosen Hilbert space?
\end{quote}

The programme supports a conditional, regime-restricted positive answer: Schr\"odinger-type evolution emerges as the formal leading-order equilibrium-tangent limit of a nonlinear closure-paired dynamics. The emergence is exact on the equilibrium solution itself and formal (first-order) near it; away from equilibrium, the dynamics is genuinely nonlinear, and full quantum mechanics --- measurement theory, Hilbert-space axioms, and the Born rule --- is \emph{not} derived here.

\subsection*{Notation}

The symbol $\sigma$ is used in the inherited Papers~XVI--XX exclusively for the Route~A Langevin noise scale; we keep that usage throughout the present paper. Paper~XXXVI's Galois/lattice automorphism of F5, which is canonically also called $\sigma$ in the cascade literature, is written here as $\sigmaG$ (Galois $\sigma$) wherever it would otherwise collide with the noise scale. The symbol $\twoI$ denotes the binary icosahedral group, a discrete subgroup of $\mathrm{SU}(2)$ of order 120; it acts on the 600-cell vertex set by isometries (the full isometry group of the 600-cell is $H_4$ of order 14400, of which $\twoI$ is a subgroup via the standard embedding). No other notation changes are introduced relative to Papers~XVI--XX and Paper~XXXVI.

%==================================================================
\section{The Five-Stage Construction}\label{sec:stages}
%==================================================================

The programme proceeds through five stages, each building on the previous:

\subsection{Stage 1: Closure evolution equation (Paper~XVI)}

Starting from the Paper~XIV~\citep{SmartXIV} Langevin equation $d\psi = -\nabla\Fcal\,dt + \sigma\,dW$ and the Paper~XV~\citep{SmartXV} complexification ansatz, a differential evolution equation is derived:
\begin{equation}\label{eq:XVI}
    \partial_t\Psi = \frac{\sigma^2}{2}\Delta\Psi - \nabla\Fcal\cdot\nabla\Psi + \frac{i}{\sigma^2}\Fcal\,\Psi
\end{equation}

This equation has the operator structure of a Fokker--Planck--Schr\"odinger hybrid: the backward Kolmogorov generator of the stochastic closure dynamics plus an imaginary (phase-generating) potential. Crucially, the kinetic term $(\sigma^2/2)\Delta$ is \textit{real diffusion}, not the imaginary propagation required for quantum evolution.

\subsection{Stage 2: Dissipative obstruction (Paper~XVII)}

Paper~XVII proves that the kinetic-term gap identified in Stage~1 is structural, not accidental:

\begin{quote}
\textit{Any generator of the form $L_{\mathrm{BK}} + iM$, where $L_{\mathrm{BK}}$ is a backward Kolmogorov generator from a real Langevin process, is dissipative in the equilibrium-weighted inner product --- regardless of the choice of imaginary potential $iM$.}
\end{quote}

This is the dissipative obstruction theorem. It establishes that phase complexification alone (Paper~XV) cannot produce unitary evolution: the obstruction sits in the kinetic sector, not the phase sector. Future progress must target a reversible kinetic construction.

\subsection{Stage 3: Nelson pairing (Paper~XVIII)}

Paper~XVIII constructs a paired forward/backward system using the closure Langevin process and its time-reversal. Defining the current velocity $v = (b_+ + b_-)/2$ and the osmotic velocity $u = (b_+ - b_-)/2$, the paper derives:
\begin{itemize}
    \item an exact continuity equation $\partial_t\rho + \nabla\cdot(\rho v) = 0$,
    \item an automatically irrotational phase function $S = -\Fcal - \sigma^2\log\sqrt{\rho}$,
    \item an exact closure Madelung pair,
    \item a nonlinear wave equation with the correct imaginary kinetic structure:
\end{itemize}
\begin{equation}\label{eq:XVIII}
    i\sigma^2\partial_t\Psi = -\frac{\sigma^4}{2}\Delta\Psi + \left(-\sigma^2\VW + \sigma^4\frac{\Delta|\Psi|}{|\Psi|}\right)\Psi
\end{equation}

The imaginary kinetic term $-(\sigma^4/2)\Delta$, absent in the single-generator framework, emerges from the forward/backward pairing. This evades the XVII obstruction because the paired dynamics is not of the $L_{\mathrm{BK}} + iM$ form.

\subsection{Stage 4: Residual isolation (Paper~XIX)}

Paper~XIX rewrites the dynamics in equilibrium-centred normal form:
\begin{equation}\label{eq:XIX}
    i\,\partial_t\Psi = \HW\Psi + \dUrel[\Psi]\,\Psi
\end{equation}
where $\HW = -(\sigma^2/2)\Delta + \VW$ is the Witten Hamiltonian and:
\begin{equation}\label{eq:residual}
    \dUrel[\Psi] = \sigma^2\left(\frac{\Delta|\Psi|}{|\Psi|} - \frac{\Delta\Rst}{\Rst}\right)
\end{equation}

The residual is proved (in Paper~XIX) to be: phase-invariant, modulus-dependent, local away from nodal sets, exactly zero at equilibrium, and $O(\varepsilon)$ near equilibrium. Amplitude homogeneity of the residual structure is proved in Paper~XX~\citep{SmartXX} rather than XIX. The full nonlinear equation conserves the $L^2$ norm exactly.

\subsection{Stage 5: Structural classification (Paper~XX)}

Paper~XX identifies the closure equation as a derived modulus-curvature nonlinear Schr\"odinger structure:
\begin{itemize}
    \item the nonlinear term belongs to the same broad class as Doebner--Goldin-type modulus-curvature corrections,
    \item the coefficient $\sigma^2$ (in Schr\"odinger-normalised form) and the equilibrium subtraction $2\VW$ are fixed by the closure geometry, not introduced phenomenologically,
    \item the linear Schr\"odinger equation arises as the leading-order (formal) linearisation at the equilibrium state, with the nonlinear residual vanishing exactly there and $O(\varepsilon)$ near it.
\end{itemize}

%==================================================================
\section{Route B: The H${}_4$-Cascade Direct Derivation}\label{sec:route_b}
%==================================================================

The stochastic route of Section~\ref{sec:stages} arrives at the Schr\"odinger equation through five stages of dynamical construction. The cascade framework of Paper~XXXVI~\citep{SmartXXXVI} supplies a conditional operator-theoretic route once the H${}_4$ generator is posited: \emph{given} the modelling assumption that evolution on the H${}_4$ rung is generated by the H${}_4$ graph Laplacian, the Schr\"odinger equation follows by Stone's theorem applied to a self-adjoint generator on a finite-dimensional Hilbert space. This route is conditional on the generator identification; it does not by itself derive the form of the Hamiltonian, which is the content of Route~A. The point of stating Route~B separately is to exhibit that once a self-adjoint operator on the H${}_4$ rung is in hand, the Schr\"odinger-type evolution follows from standard operator-theoretic results, not from any further dynamical postulate beyond the chosen generator and the standard complex Hilbert-space / operator-theory framework already in force.

\subsection{Setup: time, state space, and generator assumption}

\paragraph{Time parameter.} Paper~XXV establishes a chain-length separation on comparable pairs of H${}_4$ events and a graph-metric structure on the H${}_4$ event poset; it does not itself supply a global real-valued time parameter. Route~B adopts, as a modelling step, a real parameterisation $t \in \mathbb{R}$ of the cascade's temporal direction; this parameter is not supplied by Paper~XXV or by Paper~XXXVI F6 (F6 being a statement about spatial Gromov--Hausdorff convergence of the refined cascade lattice, not about temporal parameterisation). The identification of this parameter with standard physical time is a modelling assumption, not a theorem.

\paragraph{State space.} The Hilbert space of states at a given H${}_4$ event is the space of square-summable complex functions on the 600-cell vertex set:
\begin{equation}
    \mathcal{H} \;=\; L^2(V_{600}) \;\cong\; \mathbb{C}^{120}.
\end{equation}
Any continuum embedding of $\mathcal{H}$ into $L^2(S^3)$ would require spectral results beyond Paper~XXXVI F6 (which gives only metric Gromov--Hausdorff convergence, not operator- or spectral-convergence); no such embedding is used in the present section, and no such theorem is asserted here.

\paragraph{Generator assumption.} Route~B posits that the generator of time evolution on $\mathcal{H}$ is the H${}_4$ graph Laplacian
\begin{equation}\label{eq:H4_laplacian}
    (\Delta_{H_4}\psi)(v) \;=\; \sum_{w \sim v}\bigl[\psi(v) - \psi(w)\bigr],
\end{equation}
where $v \sim w$ denotes adjacency in the 600-cell edge graph (each vertex has 12 neighbours). This identification of the generator is a modelling choice: it is motivated by the geometry of the H${}_4$ rung (the 600-cell graph Laplacian is the canonical H${}_4$-symmetric discrete operator), but it is not derived from Route~A or from the cascade foundations F1--F8 of Paper~XXXVI. The rest of Route~B establishes what follows from this generator assumption.

\subsection{Self-adjointness of the chosen generator}\label{subsec:route_b_selfadjoint}

\begin{theorem}[Self-adjointness of the Route~B Hamiltonian $\hat H = \hbar\,\Delta_{H_4}$]\label{thm:cascade_selfadjoint}
For self-adjointness it suffices that $\hbar \in \mathbb{R}$; we take $\hbar > 0$ for the physical case. Then the operator $\hat H = \hbar\,\Delta_{H_4}$ on $\mathcal{H} = L^2(V_{600}) \cong \mathbb{C}^{120}$ is self-adjoint.
\end{theorem}

\begin{proof}
$\Delta_{H_4}$ acts on a finite-dimensional complex Hilbert space as a real symmetric matrix with respect to the standard basis of indicator functions on 600-cell vertices. Real symmetric matrices on finite-dimensional complex inner-product spaces are Hermitian (self-adjoint). Multiplication by a real scalar $\hbar$ preserves Hermiticity; hence $\hat H = \hbar\,\Delta_{H_4}$ is self-adjoint. \qed
\end{proof}

\begin{remark}[Spectrum of $\Delta_{H_4}$]\label{rem:H4_spectrum}
Direct computation (Paper~V, Theorem ``Eigenvalue structure'') gives the spectrum of $\Delta_{H_4} = 12I - A$ (with $A$ the 600-cell adjacency matrix) as
\[
    \{0,\; 9-3\sqrt{5},\; 12-4\varphi,\; 9,\; 12,\; 14,\; 4 + 4\varphi^2,\; 15,\; 9+3\sqrt{5}\},
\]
with multiplicities $\{1, 4, 9, 16, 25, 36, 9, 16, 4\}$ summing to $120$. The multiplicity pattern consists of perfect squares; a representation-theoretic explanation is not given here. Relating this finite-dimensional spectrum to the $\ell(\ell+2)$ Laplace--Beltrami spectrum of $S^3$ would require a spectral-convergence theorem going beyond Paper~XXXVI F6 (which supplies only metric Gromov--Hausdorff convergence); no continuum spectral-convergence claim is made in the present section.
\end{remark}

\begin{remark}[Role of F5 $\sigmaG$-invariance]\label{rem:F5_role}
Paper~XXXVI Foundation~F5 establishes $\sigmaG$-intertwining of the scalar closure functional $F = \alpha R + \beta E - \gamma Q$, under the hypotheses stated there on the coefficient structure (rational coefficients and $\mathbb{Q}(\varphi)$-defined ambient data); here $\sigmaG: \sqrt{5} \mapsto -\sqrt{5}$ is the Galois/lattice automorphism of the icosian construction. This is a statement about the functional $F$, not about the chosen generator $\Delta_{H_4}$ or the operator $\hat H$. In the Route~B construction, F5 supplies a discrete symmetry constraint on the cascade but is \emph{not} used to establish self-adjointness or unitarity; those follow from Theorem~\ref{thm:cascade_selfadjoint} and Stone's theorem directly. It is expected, but not proved here, that the lattice $\sigmaG$-automorphism acts on $\mathcal{H}$ as a unitary commuting with $\Delta_{H_4}$ (both inherit symmetries of the 600-cell); that would make $\sigmaG$ a conserved symmetry of the Route~B dynamics, but we do not use this below.
\end{remark}

\subsection{Stone's theorem: unitary group and Schr\"odinger equation}\label{subsec:stone}

Given the self-adjoint operator $\hat H = \hbar\,\Delta_{H_4}$ from Theorem~\ref{thm:cascade_selfadjoint}, Stone's theorem supplies a one-parameter unitary group and hence the Schr\"odinger equation. The direction of the argument is: \emph{self-adjoint generator $\Rightarrow$ unitary group $\Rightarrow$ Schr\"odinger evolution}, not the reverse.

\begin{theorem}[Stone~1932, canonical form]\label{thm:stone}
Let $B$ be a self-adjoint operator (densely defined) on a complex Hilbert space $\mathcal{H}$. Then $\{U(t) = e^{-itB}\}_{t \in \mathbb{R}}$ is a strongly continuous one-parameter group of unitary operators, and $B$ is its unique self-adjoint generator (in the sense $\tfrac{1}{-i}\frac{d}{dt}U(t)\big|_{t=0} = B$ on its dense domain). Conversely, every strongly continuous one-parameter unitary group has such a generator.
\end{theorem}

(For a proof see \citep{ReedSimon1980}, Thm.~VIII.7--VIII.8.)

\begin{corollary}[Schr\"odinger evolution from the Route~B generator]\label{cor:schrodinger_from_cascade}
Fix $\hbar > 0$. Adopt the Route~B modelling postulate that the physical Hamiltonian on the H${}_4$ rung is $\hat H = \hbar\,\Delta_{H_4}$; this postulate is \emph{not} supplied by Theorem~\ref{thm:cascade_selfadjoint}, which only proves self-adjointness of that chosen operator once the postulate is in force. Set $B := \hat H / \hbar = \Delta_{H_4}$ (self-adjoint as a real scalar multiple of a self-adjoint operator). Applying Theorem~\ref{thm:stone} to $B$ produces a strongly continuous one-parameter unitary group $U(t) = e^{-itB} = e^{-it\hat H/\hbar}$ on $\mathcal{H}$, and every trajectory $|\psi(t)\rangle = U(t)|\psi(0)\rangle$ satisfies
\begin{equation}\label{eq:cascade_schrodinger}
    i\hbar\,\partial_t |\psi(t)\rangle \;=\; \hat H\,|\psi(t)\rangle, \qquad \hat H = \hbar\,\Delta_{H_4}.
\end{equation}
\end{corollary}

\begin{proof}
By Theorem~\ref{thm:cascade_selfadjoint}, $\hat H$ and hence $B = \hat H/\hbar$ is self-adjoint on $\mathcal{H}$. Applying Theorem~\ref{thm:stone} with this $B$ produces the unitary group $U(t) = e^{-itB}$; substituting $B = \hat H/\hbar$ gives $U(t) = e^{-it\hat H/\hbar}$. Differentiating $|\psi(t)\rangle = U(t)|\psi(0)\rangle$ in $t$ and multiplying by $i\hbar$ yields~\eqref{eq:cascade_schrodinger}. \qed
\end{proof}

\begin{remark}[What is and is not derived here]\label{rem:route_b_scope}
The chain of Theorem~\ref{thm:cascade_selfadjoint} and Corollary~\ref{cor:schrodinger_from_cascade} establishes that \emph{if} the cascade dynamics on the H${}_4$ rung is generated by $\hbar\,\Delta_{H_4}$, \emph{then} Schr\"odinger-type evolution with that Hamiltonian follows from standard operator theory. It does \emph{not} derive the choice of generator from deeper cascade principles. The generator choice is the substantive modelling input of Route~B; the operator-theoretic conclusion from that input is essentially routine.
\end{remark}

\subsection{Continuum limit (sketch)}

Connecting equation~\eqref{eq:cascade_schrodinger} to the textbook continuum Schr\"odinger equation $i\hbar\,\partial_t\psi = [-(\hbar^2/2m)\nabla^2 + V(\mathbf{x})]\psi$ requires three separate steps, none of which is carried out in the present section and each of which is conditional on hypotheses stated in the cited paper:

\begin{enumerate}
    \item \emph{Spectral convergence.} Paper~XXXVI Foundation~F6 gives Burago--Ivanov Gromov--Hausdorff convergence of the refined cascade lattice to $S^3$ with a provisional covering-radius rate; it does not itself state a spectral-convergence theorem for $\Delta_{H_4}$. Deriving convergence of the $\Delta_{H_4}$ spectrum, restricted to the lowest $N$ eigenspaces, to the Laplace--Beltrami spectrum $\{\ell(\ell+2)\}$ of $S^3$ would require a further theorem (e.g., of Cheeger--Colding type~\citep{CheegerColding1997}) that is not proved in the cited source.
    \item \emph{Flat-space limit.} On length scales small compared to the $S^3$ radius, $-\Delta_{S^3} \to -\nabla^2_{\text{Euclidean}}$ at leading order; this is standard Riemannian-geometry asymptotics and is not reproved here.
    \item \emph{Potential term.} A multiplication operator $V(\mathbf{x})$ on the continuum Hilbert space must be constructed from the cascade functional's $\alpha R$ component (or from an external coupling, for test-particle problems). The systematic derivation $\alpha R \rightsquigarrow V(\mathbf{x})$ is not carried out in the present paper and remains an open step in the Route~B programme.
\end{enumerate}

Under (1)--(3), one obtains a conditional route from the H${}_4$ generator model of Corollary~\ref{cor:schrodinger_from_cascade} to the textbook continuum Schr\"odinger equation. None of the three steps is carried out here; each rests on results or open questions in the papers cited above. The effective parameters ($\hbar$ as cascade action scale, $m$ as shell-depth mass) would be determined by the cascade foundations only up to the normalisations fixed conditionally in Paper~XXXVI F1+F8 (see Paper~XXXVI's explicit caveats about imported couplings in the F8 subsection). The Route~B programme therefore trades the Schr\"odinger equation as a primitive postulate for the generator choice of Section~\ref{subsec:route_b_selfadjoint}; this is a relocation of the axiomatic input, not its elimination. It does not remove all free parameters, and a full parameter-counting analysis is outside the scope of the present paper.

\subsection{Routes~A and~B as parallel structural paths}\label{subsec:routes_parallel}

Routes~A and~B both produce Schr\"odinger-type equations $i\hbar\,\partial_t\psi = \hat H\psi$, but with \emph{different} Hamiltonians and under different supporting assumptions:

\begin{itemize}
    \item Route~A: $\hat H = \HW = -(\sigma^2/2)\Delta + \VW$ (Witten Hamiltonian from closure dynamics, via Nelson pairing). Valid at the equilibrium tangent of a nonlinear dynamics; $\sigma^2$ is a noise scale.
    \item Route~B: $\hat H = \hbar\,\Delta_{H_4}$ (H${}_4$ graph Laplacian). Valid under the generator assumption of Section~\ref{subsec:route_b_selfadjoint}; $\hbar$ is the cascade action scale.
\end{itemize}

The two Hamiltonians are not presently identified. Matching $\HW$ to $\hbar\,\Delta_{H_4}$ at the level of effective parameters requires both (i) an operator/spectral continuum limit from $\Delta_{H_4}$ to $-\Delta_{S^3}$ (which would require a theorem extending F6 from Gromov--Hausdorff convergence to operator/spectral convergence, e.g., of Cheeger--Colding~\citep{CheegerColding1997} type; F6 itself does not supply this) and (ii) an identification $\VW \leftrightarrow V(\mathbf{x})$ from an $\alpha R$-coupling analysis (Section~\ref{sec:route_b}, step~(3) of the continuum-limit sketch). Neither step is carried out here.

What Routes~A and~B share, in their \emph{linear regimes}, is evolution governed by self-adjoint operators on complex Hilbert spaces: Route~A supplies such an operator as the formal equilibrium-tangent linearisation of an exact nonlinear dynamics on the full state space, under the finite-dimensional invariant-subspace hypothesis of Section~\ref{sec:geometric}; Route~B supplies one directly as the posited generator on the finite-dimensional H${}_4$ rung. The two levels of statement are not the same: Route~B is a conditional operator-theory exercise once the generator is posited; Route~A is a formal equilibrium-tangent recovery from a broader nonlinear dynamical construction. The claim is \emph{not} that Routes~A and~B are two derivations of the same equation; they are two parallel paths, at different levels of evidential strength, that land in the same equation-class.

%==================================================================
\section{The Central Structural Result}\label{sec:result}
%==================================================================

The combined outcome of Papers~XVI--XX is:

\begin{quote}
\textbf{Structural synthesis.} Papers~XVIII--XX establish an exact nonlinear paired-process wave equation whose equilibrium state has vanishing residual, and whose leading-order linearisation at that equilibrium coincides with the Schr\"odinger equation with Witten Hamiltonian $\HW$. The linearisation is formal/perturbative in the small parameter $\varepsilon$ measuring deviation from equilibrium; remainder estimates and rigorous regularity hypotheses are beyond the scope of XVIII--XX and of this synthesis.
\end{quote}

More precisely (in Route~A's natural units, where the reduced action scale is set to unity so that $\hbar$-dependence is carried by the noise scale $\sigma^2$; equivalently, Route~A takes $\hbar_{\mathrm{eff}} = 1$ and uses $\sigma^2$ as its kinetic coefficient):
\begin{itemize}
    \item The full dynamics is an exact nonlinear Schr\"odinger-type equation with a state-dependent modulus-curvature term.
    \item The nonlinear term vanishes exactly at the equilibrium amplitude, and more generally when the amplitude has the same curvature-to-amplitude ratio $\Delta|\Psi|/|\Psi|$ as the equilibrium profile (cf.~equation~\eqref{eq:residual}).
    \item Near equilibrium, the dynamics reduces to:
    \begin{equation}
        i\,\partial_t\Psi = \HW\Psi = \left(-\frac{\sigma^2}{2}\Delta + \VW\right)\Psi
    \end{equation}
    \item The effective parameters ($m_{\mathrm{eff}} = 1/\sigma^2$, $V_{\mathrm{eff}} = \VW$) are identified in the equilibrium-tangent normal form by the closure functional $\Fcal$ and noise scale $\sigma$; this is a formal identification inside Route~A's natural units, not a physical determination in standard QM units.
\end{itemize}

Restoring $\hbar$ in physical units suggests the structural correspondence $\hbar_{\mathrm{Route\,A}} \leftrightarrow \sigma^2$ between the Route~A noise scale and Planck's constant; this is a bookkeeping correspondence at the level of units, not an established physical identification (the physical identification is explicitly disclaimed in Section~\ref{sec:not}). The recovery-map entries for $\hbar$ (Section~\ref{sec:recovery}) record the correspondence in this conditional sense. Within this framework, the Schr\"odinger equation appears formally as the leading-order equilibrium-tangent sector of a broader nonlinear dynamics, rather than as a starting axiom.

\begin{remark}[On non-triviality of the tangent limit]
While any nonlinear system admits a linearisation at a fixed point, the non-trivial content here is that the linearised generator is explicitly the Witten Hamiltonian $\HW$, derived from the closure functional $\Fcal$ and noise scale $\sigma$, rather than an arbitrary linear operator.
\end{remark}

%==================================================================
\section{Recovery Map to Standard Quantum Mechanics}\label{sec:recovery}
%==================================================================

The Recovery Manifest voice of the VFD programme holds that standard physics is not wrong but incomplete: the cascade is proposed to supply structural counterparts for standard-QM objects, with some of these being theorems of the present paper, others being postulates of the chosen framework, and others being conjectural identifications. Table~\ref{tab:recovery_qm} records the explicit correspondence, with the epistemic status of each entry indicated in its description.

\begin{table}[ht]
\centering
\caption{Recovery map: standard-QM objects and their cascade-side counterparts or modelling inputs.}
\label{tab:recovery_qm}
\begin{tabular}{@{}p{5cm}p{8cm}@{}}
\toprule
Standard-QM object & Cascade-side counterpart or modelling input \\
\midrule
State $|\psi\rangle$ (wavefunction) & Complex function on 600-cell vertices. A possible continuum embedding into $L^2(S^3)$ would require spectral results beyond F6 and is not used here \\
Hilbert-space inner product & Standard $L^2$ form on 600-cell; $\twoI$ acts by isometries. A unitary $\sigmaG$ action on $\mathcal{H}$ is conjectural here (Remark~\ref{rem:F5_role}) \\
Time parameter $t$ & Paper~XXV event-order structure provides a chain-length separation; a global real-valued $t$ is adopted as a modelling parameterisation of the cascade's temporal direction. It is not supplied by Paper~XXV or by Paper~XXXVI F6 \\
Hamiltonian $\hat H$ & Route~A: Witten Hamiltonian $\HW$ from closure functional at equilibrium tangent. Route~B: posited generator $\hbar\,\Delta_{H_4}$ (H${}_4$ graph Laplacian) \\
Evolution operator $U(t)$ & Route~A: formal equilibrium-tangent unitary evolution by $\HW$ on the invariant finite-dimensional sector subspace (Hypothesis~\ref{hyp:subspace_invariance}). Route~B: $U(t) = e^{-it\hat H/\hbar}$ unitary by Stone's theorem applied to Theorem~\ref{thm:cascade_selfadjoint} \\
Planck constant $\hbar$ & Route~A: bookkeeping correspondence with the noise scale $\sigma^2$, not a physical identification. Route~B: cascade action scale, conditional on Paper~XXXVI F1+F8 normalisations \\
Schr\"odinger equation & Route~A: formal equilibrium-tangent linearisation with $\HW$ (Section~\ref{sec:result}). Route~B: Stone's theorem for the posited self-adjoint H${}_4$ graph Laplacian (Corollary~\ref{cor:schrodinger_from_cascade}) \\
Self-adjointness of the chosen Hamiltonian/generator & Route~B: real-symmetric structure of the H${}_4$ graph Laplacian gives self-adjointness on $\mathcal{H}$ (Theorem~\ref{thm:cascade_selfadjoint}). No general ``self-adjointness of observables'' theorem is claimed here \\
Spectral discreteness & 9 distinct H${}_4$ Laplacian eigenvalues (direct computation; see Remark~\ref{rem:H4_spectrum}) \\
Degeneracy structure & Perfect-square multiplicities, suggestive of a $\twoI$-representation-theoretic explanation; no precise decomposition is proved in the present paper \\
\bottomrule
\end{tabular}
\end{table}

The recovery map does not rename standard objects; it identifies each with a cascade-combinatorial counterpart that carries additional structure. In Route~B, once the generator $\hat H = \hbar\,\Delta_{H_4}$ is posited, its self-adjointness on the finite-dimensional $\mathcal{H}$ is immediate (Theorem~\ref{thm:cascade_selfadjoint}), and Schr\"odinger evolution follows by Stone's theorem applied to that self-adjoint generator (Corollary~\ref{cor:schrodinger_from_cascade}). In Route~A, self-adjointness of the effective linearised Hamiltonian is used only on the finite-dimensional invariant subspace of Hypothesis~\ref{hyp:subspace_invariance}; operator/domain hypotheses for $\HW$ on the full state space are not established here. In both routes, once a self-adjoint generator is in hand in its declared regime, Schr\"odinger evolution is an operator-theoretic consequence rather than an additional operator-theoretic postulate.

\paragraph{What this means for QM's status.} Standard quantum mechanics remains empirically correct within its tested regime. The cascade framework adds structural \emph{supplements}, not a replacement:
\begin{enumerate}[label=(\alph*)]
    \item an account of how Schr\"odinger evolution can arise from a generator choice plus standard operator theory, rather than being a free postulate;
    \item a specific computation of the H${}_4$ graph Laplacian spectrum whose numerical values are used as inputs to shell-depth mass correspondences in Paper~V (Paper~V presents these as model-level numerical agreements, not first-principles theorems);
    \item a proposed structural correspondence for $\hbar$ (Route~A: bookkeeping correspondence with the noise scale $\sigma^2$, not a physical identification; Route~B: cascade action scale, conditional on F1+F8 normalisations in Paper~XXXVI);
    \item a predicted nonlinear residual (Route~A) at modulus-curvature scale $O(\varepsilon)$ that may or may not be experimentally accessible.
\end{enumerate}
Measurement theory and the Born rule are \emph{not} derived in the present paper; they are the subject of Papers~XXX and~XXXI. Point~(a) supplies a structural supplement to standard QM's axiomatic starting point; points~(b)--(c) supply prospective structural identifications for quantities treated as free parameters in standard QM; point~(d) is a cascade-specific extension testable in principle. None of (a)--(d) claims to complete standard QM on its own.

%==================================================================
\section{Geometric Structure at the Equilibrium Tangent}\label{sec:geometric}
%==================================================================

\begin{remark}[Addendum status and scope]\label{rem:addendum_geometric}
This section was added in a subsequent revision to make explicit the geometric-quantum structure of the equilibrium tangent regime identified in Section~\ref{sec:result}. The disclaimers of Section~\ref{sec:not} are unchanged: unitarity and projective Hilbert structure are \emph{not} claimed for the full nonlinear dynamics of~\eqref{eq:XVIII}--\eqref{eq:XIX}. The content here is restricted to the leading-order linearisation at equilibrium, $i\,\partial_t\Psi = \HW\Psi$ (in the perturbative sense of Papers~XVIII--XX), further restricted to a finite-dimensional sector-superposition subspace under the explicit Hypothesis~\ref{hyp:subspace_invariance} below (sector-subspace invariance under $\HW$). The projective/K\"ahler geometry is standard geometric quantum mechanics applied to this finite-dimensional subspace; the closure-derived content is the identification of $\HW$ as the effective Hamiltonian on that subspace, not the projective geometry itself. Within the stated scope, the section assembles the Riemannian, symplectic, complex, and unitary structures that the downstream Born-rule derivation of Paper~XXX uses in its K\"ahler argument.
\end{remark}

At the equilibrium tangent regime, the dynamics is formally generated by the Witten Hamiltonian $\HW$ acting on a linear space of perturbations around the equilibrium state $\Psi_{\mathrm{st}}$. For the geometric discussion below we further restrict to a finite-dimensional subspace on which the restricted operator is assumed self-adjoint. We state the invariance and self-adjointness requirements as an explicit hypothesis:

\begin{definition}[Finite sector-superposition subspace]\label{def:sector_subspace}
Let $\mathcal{H}_{\mathrm{amb}}$ denote the ambient Hilbert space in which the equilibrium-tangent linearisation is realised; in the present section its precise construction is not needed, only its inner product. Let $\mathcal{H}_O \subset \mathcal{H}_{\mathrm{amb}}$ denote a finite-dimensional complex subspace of dimension $N \geq 2$, spanned by sector basis vectors $\phi_1, \ldots, \phi_N$ associated with chosen basins of $\Fcal$. $\mathcal{H}_O$ inherits the inner product $\langle\cdot,\cdot\rangle$ from $\mathcal{H}_{\mathrm{amb}}$.
\end{definition}

\begin{hypothesis}[Sector-subspace invariance and restricted self-adjointness]\label{hyp:subspace_invariance}
$\mathcal{H}_O$ is invariant under the equilibrium-tangent Witten Hamiltonian: $\HW \mathcal{H}_O \subseteq \mathcal{H}_O$. We further assume that the restricted operator $\HW|_{\mathcal{H}_O}$ is self-adjoint on the finite-dimensional Hilbert space $\mathcal{H}_O$. (Invariance alone is not enough to imply self-adjointness of the restriction: one also needs a symmetry of $\HW$ on a suitable dense domain. We take both requirements as part of the hypothesis.)
\end{hypothesis}

\begin{remark}[Status of Hypothesis~\ref{hyp:subspace_invariance}]\label{rem:hyp_status}
This is a genuine assumption, not a consequence of closure dynamics. A sufficient condition is that $\mathcal{H}_O$ is spanned by eigenvectors of $\HW$; in general, whether a chosen sector decomposition yields an $\HW$-invariant subspace depends on the relation between the basin structure and the spectral structure of $\HW$. It is taken here as the defining condition of the regime in which the geometric-quantum identification applies. Paper~XXX adopts the observer-side completeness hypothesis $\Mcl(\Ocal) = \mathbb{CP}^{N-1}$, which is a compatible but logically distinct assumption: completeness is about which states are observer-admissible, whereas Hypothesis~\ref{hyp:subspace_invariance} is about which subspace is preserved by $\HW$. Paper~XXX's argument needs both.
\end{remark}

\subsection{Projective state manifold}\label{subsec:cpN}

Pure states of the equilibrium-tangent dynamics, restricted to $\mathcal{H}_O$, are unit vectors in $\mathcal{H}_O$ modulo global phase. Writing $S(\mathcal{H}_O) = \{\Psi \in \mathcal{H}_O : \|\Psi\| = 1\}$ for the unit sphere and $U(1)$ acting by $\Psi \mapsto e^{i\alpha}\Psi$, the projective state manifold is the quotient
\begin{equation}\label{eq:CPN}
    \mathbb{CP}^{N-1} \;=\; S(\mathcal{H}_O) / U(1).
\end{equation}
A point $[\Psi] \in \mathbb{CP}^{N-1}$ represents the $U(1)$-orbit of a unit vector $\Psi \in \mathcal{H}_O$. The finite-dimensional projective structure is well-defined precisely because $\mathcal{H}_O$ is finite-dimensional (Definition~\ref{def:sector_subspace}); an infinite-dimensional ambient Hilbert space would require a more delicate construction.

\subsection{Hermitian, Riemannian, and symplectic structures}\label{subsec:structures}

On $\mathcal{H}_O$ the Hermitian inner product $h(\Psi, \Psi') = \langle\Psi|\Psi'\rangle$ decomposes into real and imaginary parts, which at a unit vector $\Psi$ define a Riemannian metric $g$ and a symplectic form $\omega$ on the tangent space:
\begin{align}
    g(v, w) &\;=\; \mathrm{Re}\,\langle v|w\rangle \qquad && \text{(Riemannian, on }T_\Psi \mathcal{H}_O\text{)}, \label{eq:g_from_h} \\
    \omega(v, w) &\;=\; \mathrm{Im}\,\langle v|w\rangle \qquad && \text{(symplectic, on }T_\Psi \mathcal{H}_O\text{)}, \label{eq:omega_from_h}
\end{align}
where $v, w \in T_\Psi \mathcal{H}_O \cong \mathcal{H}_O$. These are ambient forms; to obtain forms on $\mathbb{CP}^{N-1}$ we restrict to the unit sphere $S(\mathcal{H}_O)$ and quotient by the global phase. Explicitly: the tangent space $T_{[\Psi]}\mathbb{CP}^{N-1}$ is identified with the horizontal subspace $\{v \in T_\Psi S(\mathcal{H}_O) : \langle\Psi|v\rangle = 0\}$ (orthogonal to both radial and phase directions), and the restrictions of $g$ and $\omega$ to this subspace descend to the quotient. The resulting forms on $\mathbb{CP}^{N-1}$ are, by the standard geometric-quantum construction (see \citep{AshtekarSchilling1999, KobayashiNomizu1969}), a constant multiple of the Fubini--Study metric and its associated K\"ahler form; the scale depends on the chosen normalisation of $h$.

The almost-complex structure $J$ defined by multiplication by $i$ (see Paper~XVIII~\citep{SmartXVIII}, Section on Nelson geometric structure) preserves the horizontal subspace and descends to $\mathbb{CP}^{N-1}$. It is compatible with the descended $g$ and $\omega$:
\begin{equation}\label{eq:compat}
    g(v, w) \;=\; \omega(v, J w) \qquad \text{on horizontal tangents},
\end{equation}
which is the K\"ahler compatibility relation. The triple $(g, J, \omega)$ on $\mathbb{CP}^{N-1}$ is a K\"ahler structure; it coincides with the standard Fubini--Study K\"ahler structure up to the scale factor induced by the normalisation of $h$.

\subsection{Equilibrium-tangent Schr\"odinger dynamics is Hamiltonian}\label{subsec:hamiltonian}

\begin{proposition}[Formal equilibrium-tangent linearised Schr\"odinger flow on $\mathbb{CP}^{N-1}$; imported standard fact of geometric QM \citep{AshtekarSchilling1999}]\label{prop:schroedinger_hamiltonian}
Under Hypothesis~\ref{hyp:subspace_invariance}, the formal equilibrium-tangent linearised Schr\"odinger flow $i\,\partial_t\Psi = \HW\Psi$ of Stage~5, restricted to $\mathcal{H}_O$, induces the standard projective Hamiltonian flow on $\mathbb{CP}^{N-1}$ with respect to the descended symplectic form (the projection of $\omega$ from~\eqref{eq:omega_from_h} onto horizontal tangents). The Hamiltonian function is $h_{\HW}([\Psi]) = \langle\Psi|\HW|\Psi\rangle/\langle\Psi|\Psi\rangle$.
\end{proposition}

\begin{proof}
For any self-adjoint operator $A$ on $\mathcal{H}_O$, the function $h_A([\Psi]) = \langle\Psi|A|\Psi\rangle/\langle\Psi|\Psi\rangle$ is smooth on $\mathbb{CP}^{N-1}$. Its Hamiltonian vector field with respect to $\omega$ is $X_{h_A}([\Psi]) = -i(A - h_A \mathbb{1})\Psi$ (a standard computation; see~\citep{AshtekarSchilling1999}). Taking $A = \HW$ and writing the flow of $X_{h_{\HW}}$ at the base point:
\begin{equation}
    \dot\Psi \;=\; -i\,\HW\,\Psi \;-\; \text{(projective phase term)},
\end{equation}
which reduces to $i\,\partial_t\Psi = \HW\Psi$ after passing to $\mathbb{CP}^{N-1}$ (the projective phase term reflects only the quotient by global phase). Therefore the Schr\"odinger dynamics at the equilibrium tangent is the Hamiltonian flow of the expectation-value function of $\HW$.
\end{proof}

\subsection{Unitary action and $U(N)$ covariance}\label{subsec:unitary}

\begin{proposition}[Induced projective action of $U(N)$ preserves the K\"ahler structure; imported standard fact \citep{KobayashiNomizu1969}]\label{prop:unitary}
Every unitary $U \in U(\mathcal{H}_O) = U(N)$ acts on $\mathbb{CP}^{N-1}$ through its induced projective action, which factors through $\mathrm{PU}(N) = U(N)/U(1)$. This projective action preserves the descended K\"ahler structure $(g, J, \omega)$ on $\mathbb{CP}^{N-1}$: $g$, $\omega$, and $J$ are all invariant under the induced map. Conversely, the holomorphic isometry group of this K\"ahler structure on $\mathbb{CP}^{N-1}$ is $\mathrm{PU}(N)$, by the standard classification of K\"ahler isometries (cited below).
\end{proposition}

\begin{proof}
By definition of unitarity, $\langle U\Psi|U\Psi'\rangle = \langle\Psi|\Psi'\rangle$ for all $\Psi, \Psi' \in \mathcal{H}_O$. Taking real and imaginary parts gives $g(U\cdot, U\cdot) = g$ and $\omega(U\cdot, U\cdot) = \omega$ as ambient forms on $T\mathcal{H}_O$; these descend to the projective quotient because $U$ preserves the unit sphere and commutes with the $U(1)$-action (global phase). Compatibility with $J$ follows because multiplication by $i$ commutes with any complex-linear map. The converse statement --- that every K\"ahler isometry of $\mathbb{CP}^{N-1}$ arises this way --- is a standard classification result, imported here as a citation~\citep{KobayashiNomizu1969} rather than reproved.
\end{proof}

\begin{corollary}[Equilibrium-tangent Schr\"odinger flows give unitary action on $\mathcal{H}_O$ and projective K\"ahler isometries]\label{cor:schroedinger_unitary}
For any self-adjoint $A$ on $\mathcal{H}_O$, the time-$t$ flow $e^{-iAt}$ is unitary on $\mathcal{H}_O$ (a one-parameter subgroup of $U(N)$) and its induced projective action on $\mathbb{CP}^{N-1}$ factors through $\mathrm{PU}(N) = U(N)/U(1)$, acting there as a K\"ahler isometry. In particular, the restricted linear Schr\"odinger flow generated by $\HW|_{\mathcal{H}_O}$ (Proposition~\ref{prop:schroedinger_hamiltonian}) sits in $U(N)$ on $\mathcal{H}_O$ and induces a one-parameter subgroup of $\mathrm{PU}(N)$ on $\mathbb{CP}^{N-1}$.
\end{corollary}

\begin{proof}
$e^{-iAt}$ is unitary on $\mathcal{H}_O$ for $A = A^\dagger$. Apply Proposition~\ref{prop:unitary}.
\end{proof}

\subsection{What is and is not claimed}\label{subsec:xxi_scope}

\begin{remark}[Scope of this section]\label{rem:scope_geometric}
Propositions~\ref{prop:schroedinger_hamiltonian} and~\ref{prop:unitary} apply strictly to the equilibrium tangent regime where the dynamics is $i\,\partial_t\Psi = \HW\Psi$. The full nonlinear closure-Nelson equation~\eqref{eq:XVIII} does \emph{not} generate a $U(N)$ flow in general, because the state-dependent residual potential breaks $U(N)$-equivariance. The equilibrium tangent is the sector in which the geometric-quantum identification holds; this is consistent with, and does not extend, the disclaimers of Section~\ref{sec:not}.
\end{remark}

\begin{remark}[Role in Paper~XXX]\label{rem:role_xxx_geo}
Paper~XXX uses the K\"ahler structure $(g, J, \omega)$ of this section, together with the closure Hessian (which supplies an alternative Riemannian metric), as one imported ingredient in a conditional Petz-plus-Gleason argument for the Born rule. That argument is not carried out here, and is not a consequence of the present paper alone: Paper~XXX adopts additional observer-side hypotheses (universal completeness $\Mcl(\Ocal) = \mathbb{CP}^{N-1}$, continuity, eigenstate certainty, restricted non-contextuality, and for $N \geq 3$ the Gleason step) under which Petz's uniqueness theorem identifies the closure-induced K\"ahler structure with Fubini--Study and Gleason's frame-function theorem delivers the Born rule as the unique compatible probability assignment. The present paper's contribution to that chain is limited to making explicit the equilibrium-tangent geometric-quantum structure, with the projective/K\"ahler geometry imported from standard geometric QM (Propositions~\ref{prop:schroedinger_hamiltonian}, \ref{prop:unitary}) and the cascade-derived content being the identification of $\HW$ as the effective Hamiltonian on the invariant subspace.
\end{remark}

%==================================================================
\section{Relation to Known Formulations}\label{sec:relation}
%==================================================================

\subsection{Madelung / Bohmian mechanics}

Madelung~\citep{Madelung1927} showed that the Schr\"odinger equation, under the polar decomposition $\psi = \sqrt\rho\,e^{iS/\hbar}$, is equivalent to a pair of fluid-dynamical equations (continuity plus Hamilton--Jacobi) with a single additional term: the \emph{quantum potential} $Q = -(\hbar^2/2m)\,\Delta\sqrt\rho/\sqrt\rho$. Bohm~\citep{Bohm1952a,Bohm1952b} developed this hydrodynamic structure into the pilot-wave / de~Broglie--Bohm formulation of QM. The closure construction of Paper~XVIII produces the same modulus-curvature functional, but with a structural difference:
\begin{itemize}
    \item In standard Madelung/Bohm, the quantum potential is absorbed entirely into the linear Schr\"odinger kinetic term.
    \item In the closure framework, a sign reversal (inherited from the dissipative origin of the forward process) prevents full absorption, leaving a nonlinear residual with an equilibrium subtraction.
\end{itemize}

The closure residual is therefore the non-cancelled remainder of the same object that, in linear QM, is hidden inside the Schr\"odinger kinetic structure. Whether $\Delta\sqrt\rho/\sqrt\rho$ admits a Route-B interpretation in terms of the H${}_4$ graph Laplacian's action on amplitude modes in a suitable continuum/amplitude-sector limit is a natural question raised by the parallel between the two routes; we do not derive such an interpretation here.

\subsection{Nelson stochastic mechanics}

Nelson~\citep{Nelson1966,Nelson1985} derived the Schr\"odinger equation from a postulated stochastic Newton equation $ma(X_t) = -\nabla V(X_t)$ combined with a Brownian-diffusion hypothesis for the particle trajectory. The pairing construction of Paper~XVIII is structurally analogous: both derive Schr\"odinger dynamics from forward/backward stochastic processes. The differences are:
\begin{itemize}
    \item Nelson postulates a stochastic Newton equation ($ma = -\nabla V$) that produces the correct quantum-potential sign.
    \item The closure framework derives the dynamics from the Fokker--Planck structure of a gradient-flow Langevin equation, which produces the reversed sign.
    \item The closure construction does not require postulating an acceleration law; it uses only the existing stochastic process and its time-reversal.
\end{itemize}

The price of derivation (rather than postulation) is the nonlinear residual. The benefit is that Route~A does not start from the standard quantum dynamical postulates; Route~B instead assumes the standard complex Hilbert-space/operator-theory framework plus a chosen self-adjoint generator. Historical critiques of Nelson's programme~\citep{Wallstrom1994} point out that the quantisation condition on the phase $S$ is postulated rather than derived. Route~B handles this differently: its unitarity comes from the self-adjoint generator postulate and Stone's theorem, and F5 contributes only a discrete symmetry constraint. The Galois-twist $\sqrt{5} \mapsto -\sqrt{5}$ action on the icosian lattice is a structural feature of the cascade independent of the unitarity question, and is not used here to derive the phase.

\subsection{Nonlinear Schr\"odinger equations}

The closure equation belongs to the broad class of modulus-curvature nonlinear Schr\"odinger equations. Within this class, the closure equation is distinguished by:
\begin{itemize}
    \item a coefficient ($\sigma^2$ in normalised form) fixed by the noise scale,
    \item an equilibrium subtraction ($2\VW$ in the Schr\"odinger-normalised form of Papers~XIX--XX) fixed by the closure geometry,
    \item a distinguished linear limit selected by the equilibrium state.
\end{itemize}

No exact identification with the full Doebner--Goldin family \citep{DoebnerGoldin1996} or any specific existing model is claimed. The closure equation is a \textit{derived member} of a recognised class, not a phenomenological addition. Experimental bounds on modulus-curvature nonlinear deviations from standard QM \citep{BialynickiBirulaMycielski1976,WeinbergNonlinear1989,BollingerHeinzenItano1989} constrain the cascade residual; translating the closure residual's $\sigma^2 (\Delta|\Psi|/|\Psi| - \Delta R_{\mathrm{st}}/R_{\mathrm{st}})$ into a testable experimental deviation is an open question (Section~\ref{sec:open}).

\subsection{Guerra--Morato and Euclidean stochastic mechanics}

Guerra and Morato~\citep{GuerraMorato1983} developed an action-principle formulation of Nelson's stochastic mechanics that recovers the Schr\"odinger equation from a variational principle on forward/backward drift fields. Zambrini~\citep{Zambrini1986} extended this to a ``Euclidean quantum mechanics'' in which the Schr\"odinger equation arises from a Bernstein-process characterisation. The closure programme's Nelson-pairing construction (Paper~XVIII) occupies a similar mathematical niche but begins from a closure functional and gradient-flow Langevin equation rather than from a postulated quantum action. The relation is structural-analogue rather than direct-equivalence: both frameworks extract Schr\"odinger dynamics from paired stochastic structures, but they differ in what supplies the input drift.

\subsection{Quantum mechanics reconstruction programmes}

Several programmes in the foundations literature derive quantum theory from information-theoretic axioms or operational principles: Hardy~\citep{Hardy2001} reconstructs finite-dimensional QM from five axioms; Masanes and M\"uller~\citep{MasanesMuller2011} reconstruct from a shorter axiom list; Chiribella, D'Ariano, and Perinotti~\citep{ChiribellaDArianoPerinotti2011} reconstruct from purification postulates; Wilce~\citep{Wilce2012} reviews the field. These programmes deliver the Hilbert-space structure but typically stop short of deriving the Schr\"odinger equation's specific form; they take unitarity of time evolution as an additional axiom. The cascade framework operates at a different level: Route~B posits the H${}_4$ graph Laplacian as generator and then obtains unitarity by standard operator theory (Stone's theorem on a self-adjoint generator on a complex Hilbert space). The two approaches are complementary: the reconstruction programmes answer ``what mathematical structure does QM have?''; the cascade proposes ``from what deeper geometric substrate might that structure be inherited?''.

%==================================================================
\section{What Has Been Established}\label{sec:established}
%==================================================================

\begin{table}[ht]
\centering
\caption{Status of results across the programme.}
\label{tab:status}
\begin{tabular}{@{}llp{6cm}@{}}
\toprule
Paper & Result & Status \\
\midrule
XVI & Closure evolution equation & Derived generator-level evolution equation \\
XVII & Dissipative obstruction & Proved (Theorem~1 of XVII) \\
XVIII & Continuity equation & Exact \\
XVIII & Irrotational phase function & Exact (structural consequence) \\
XVIII & Nonlinear wave equation & Exact \\
XVIII & Near-equilibrium Schr\"odinger & Leading order in perturbation \\
XIX & Residual structural properties & Proved (phase invariance, locality, etc.) \\
XIX & Norm conservation & Exact for full nonlinear equation \\
XIX & Gauge non-removability & Proved for phase and linear-unitary transforms \\
XX & Modulus-curvature classification & Structural placement \\
XX & Equilibrium tangent limit & Leading-order linearisation at equilibrium (formal, not a rigorous Fr\'echet theorem with remainder estimates) \\
XXI (this) & K\"ahler structure at equilibrium tangent & Assembled from standard projective-Hilbert-space geometry, conditional on Hypothesis~\ref{hyp:subspace_invariance} (sector-subspace invariance under $\HW$) \\
XXI (this) & Projective $\mathrm{PU}(N)$ action preserving the K\"ahler structure at equilibrium tangent & Standard geometric-QM fact (Proposition~\ref{prop:unitary}); $U(N)$ acts on $\mathcal{H}_O$ by unitaries, and the induced projective action on $\mathbb{CP}^{N-1}$ factors through $\mathrm{PU}(N) = U(N)/U(1)$ and is a K\"ahler isometry there, under Hypothesis~\ref{hyp:subspace_invariance} \\
XXI (this, Route B) & Self-adjointness of $\hat H = \hbar\,\Delta_{H_4}$ & Theorem~\ref{thm:cascade_selfadjoint}; real-symmetric structure of 600-cell graph Laplacian (trivial under generator assumption) \\
XXI (this, Route B) & Schr\"odinger evolution from the Route~B generator & Corollary~\ref{cor:schrodinger_from_cascade}; Stone's theorem applied to Theorem~\ref{thm:cascade_selfadjoint}. Conditional on the generator assumption of Section~\ref{subsec:route_b_selfadjoint} \\
XXI (this, Route B) & F5 as a consistency symmetry (not a unitarity source) & Remark~\ref{rem:F5_role}; a conjectural action of the lattice $\sigmaG$-automorphism on $\mathcal{H}$, commuting with $\Delta_{H_4}$, would give a conserved discrete symmetry but is not proved here. F5 itself is a statement about the scalar functional $F$ \\
\bottomrule
\end{tabular}
\end{table}

%==================================================================
\section{What Has Not Been Established}\label{sec:not}
%==================================================================

The following are \textbf{not} claimed by the programme:

\begin{itemize}
    \item \textbf{Full quantum mechanics.} No measurement theory, no Hilbert-space axioms, no Born rule as a derived physical law.
    \item \textbf{Exact global Schr\"odinger recovery.} The full dynamics is nonlinear; the Schr\"odinger equation holds only in the equilibrium-centred regime.
    \item \textbf{Unitarity of the full nonlinear dynamics.} The full nonlinear equation~\eqref{eq:XVIII} conserves norm but is not generated by a skew-adjoint operator. The equilibrium tangent limit $i\,\partial_t\Psi = \HW\Psi$ \emph{is} generated by a self-adjoint Hamiltonian and is therefore unitary on the chosen finite-dimensional sector-superposition subspace whenever Hypothesis~\ref{hyp:subspace_invariance} holds (Corollary~\ref{cor:schroedinger_unitary}); unitarity is claimed only within that restricted scope.
    \item \textbf{Physical identification of $\sigma^2$ with $\hbar$.} The structural correspondence is established; the physical identification is not.
    \item \textbf{Experimental predictions.} No observable consequences of the nonlinear residual have been derived or bounded.
    \item \textbf{Physical status of the residual.} Paper~XIX establishes gauge non-removability under phase rotations and linear-unitary transforms; whether the nonlinear correction is removable by broader classes of transformations beyond those analysed in XIX, or represents genuine beyond-QM physics or an effective-theory artefact, remains open.
    \item \textbf{No uniqueness claim.} The programme identifies one coherent route from closure dynamics to Schr\"odinger-type structure; it does not prove that this route is unique or that the forward/backward pairing is the only viable construction.
\end{itemize}

\subsection*{Limitations of the construction}

The results establish a structural pathway from closure dynamics to Schr\"odinger-type evolution, but do not constitute a derivation of full quantum mechanics. In particular: measurement theory is not derived, Hilbert-space axioms are not introduced, the Born rule is not obtained as a physical law, and the nonlinear residual is not assigned physical interpretation. The results should therefore be understood as a mathematical construction linking dissipative dynamics to quantum structure, rather than as a complete physical theory.

%==================================================================
\section{The Logical Structure}\label{sec:logic}
%==================================================================

The programme's logical chain is:
\begin{enumerate}
    \item \textbf{Gradient-flow closure dynamics} (Papers~XII--XIV): deterministic descent, stochastic noise, Boltzmann equilibrium.
    \item \textbf{Phase complexification} (Paper~XV): imaginary potential added to path integral $\rightarrow$ interference and phase structure.
    \item \textbf{Evolution equation} (Paper~XVI): differential generator derived from the complexified path integral.
    \item \textbf{Kinetic obstruction} (Paper~XVII): single-generator unitarity is impossible $\rightarrow$ reversible kinetic sector needed.
    \item \textbf{Forward/backward pairing} (Paper~XVIII): Nelson-type construction produces imaginary kinetic term $\rightarrow$ nonlinear Schr\"odinger-type equation.
    \item \textbf{Residual classification} (Papers~XIX--XX): nonlinear remainder is isolated, classified as modulus-curvature structure with equilibrium subtraction.
    \item \textbf{Equilibrium tangent limit}: formal linearisation at equilibrium yields the Schr\"odinger equation with Witten Hamiltonian to leading order.
\end{enumerate}

Each step either proves a result, isolates an obstruction, or proposes a conditional next step in response to the previous paper's unresolved issue. The chain is structurally constructive at the evidential level stated in each paper and in this paper's Sections~\ref{sec:established} and~\ref{sec:not}.

%==================================================================
\section{The Reframing}\label{sec:reframing}
%==================================================================

The programme reframes the logical status of the Schr\"odinger equation within the closure framework:

\begin{quote}
The Schr\"odinger equation is not the full dynamics. It is the \textit{distinguished equilibrium tangent sector} of an exact nonlinear closure-paired dynamics.
\end{quote}

This reframing has several consequences:
\begin{itemize}
    \item Linear quantum mechanics is not contradicted: it is recovered on the equilibrium solution and formally to leading order near equilibrium, which is the regime most directly aligned with standard linear quantum phenomenology.
    \item The nonlinear residual defines a regime beyond linear QM that may or may not be physically accessible.
    \item The Witten Hamiltonian $\HW$ --- which connects the closure functional $\Fcal$ to quantum structure via supersymmetric quantum mechanics --- is identified as the effective quantum Hamiltonian, not by postulate, but by the equilibrium tangent construction.
\end{itemize}

%==================================================================
\section{Open Questions}\label{sec:open}
%==================================================================

\begin{enumerate}
    \item Does the nonlinear residual produce observable deviations from standard QM? If so, at what energy or density scales?
    \item Can the nonlinear equation be constrained by existing experimental bounds on nonlinear quantum mechanics?
    \item Does the closure equation possess a conserved energy functional beyond the $L^2$ norm?
    \item Can the nodal divergence of the residual be regularised, or does it indicate a physical boundary condition?
    \item Is there a variational principle for the full nonlinear closure dynamics?
    \item Can the construction be extended to field-theoretic or multi-particle settings?
    \item Can the closure residual be related to existing experimentally constrained nonlinear Schr\"odinger frameworks in a quantitatively testable way?
\end{enumerate}

%==================================================================
\section{Conclusion}\label{sec:conclusion}
%==================================================================

The closure programme outlines a mathematically coherent but partly conditional pathway:
\begin{equation*}
    \text{Dissipative dynamics} \;\rightarrow\; \text{Stochastic closure} \;\rightarrow\; \text{Paired processes} \;\rightarrow\; \text{Nonlinear wave equation} \;\rightarrow\; \text{Schr\"odinger limit}
\end{equation*}

The central conclusion is:

\begin{quote}
\textit{Linear quantum evolution can arise formally as the leading-order equilibrium-tangent limit of a broader nonlinear dynamical system in Route~A, without being postulated as the starting dynamical law, and in Route~B as a consequence of a separately posited self-adjoint generator on a standard complex Hilbert space via Stone's theorem.}
\end{quote}

The significance of this result lies not in replacing quantum mechanics --- which remains the correct description of the experimentally established linear quantum regime --- but in identifying a mathematically coherent and partly conditional route by which it is modelled here as emerging from a deeper dynamical structure. In Route~A, linear quantum evolution appears as a formal equilibrium-tangent limit of a broader nonlinear system; in Route~B it follows from a separately postulated self-adjoint generator. Whether the deeper structure has physical content beyond the linear quantum limit is the defining open question for the next stage of the programme.

%==================================================================
\bibliographystyle{plainnat}
\bibliography{references}

\end{document}
